\section{Lecture 4: On measuring biodiversity and its importance}

Lars Gamfeldt

\subsection{Seminar overview}
\begin{itemize}
	\item What is biodiversity and how do we measure it?
	\item How do we know that biodiversity is declining, what metrics
		do we use?
	\item Why is it important to measure biodiversity change?
\end{itemize}

Biomass of different groups:
\begin{itemize}
	\item Plants: 450 Gt C
		\begin{itemize}
			\item Anthropods (e.g. insects): 1Gt C
			\item Fish: 0.7Gt C
			\item Wild mammals: 0.007 Gt C
			\item Humans: 0.06 Gt C
			\item Livestock: 0.1 Gt C
			\item Wild birds: 0.002 Gt C
			\item Annelids: 0.2 Gt C
		\end{itemize}
	\item Animals: 2 Gt C
	\item Fungi: 12 Gt C
	\item Bacteria: 70 Gt C
	\item Protists: 4 Gt C
	\item Viruses: 0.2 Gt C
	\item 7 Gt C
\end{itemize}

\subsection{What is biodiversity}

ecosystems $\leftrightarrow$ species $\leftrightarrow$ genes


\begin{itemize}
	\item Biodiversity is a concept not a quantity
	\item Typically, we recognise different dimensions of biodiversity
	\item We operationalise it in different ways
\end{itemize}

\subsection{Biodiversity and human health: An analogy}

physical, social and mental healths all overlap with each other

anologously, biodiversity has 4 different dimensions:
\begin{itemize}
	\item evolutionary: gene diversity, phylogenetics etc.
	\item functional: trait diversity, trophic diversity etc.
	\item taxonomic: species diversity etc.
	\item abundance: population size, biomass etc.
\end{itemize}


\begin{itemize}
	\item biodiversity consists of different dimensions
	\item use a variety of metrics to try and quantify these dimensions
	\item understanding biodiversity of a place, time etc.
\end{itemize}

\subsection{Is biodiversity higher on land or in the sea?}

Number of species on land is about 8M an about 2M in the ocean (both have
huge margin of error).

77\% of all animal species occur on land, 11\% in freshwater, 12\% are marine.
Explanation: 90\% of species of land are anthropods, most of which are insects.

\subsection{Terrestial Anthropods are extremely diverse}

Erwin (1982) surveyed insects on 19 different trees of the same species, he
found 1200 insect species.

There could be 1.5M species of beetles alone!

\subsection{Is biodiversity higher on land or in the sea?}
\begin{itemize}
	\item It depends on how it is measured
	\item Depends on the value put on different dimensions of biodiversity
\end{itemize}

\subsection{The Living Planet Index (WWF)}

It reflects the average size of wildlife populations.

\subsection{The Red List Index}

Consists of 9 categories from Extinct to Least Concern (through Threatened).

\subsection{Global extinction rates}

Current extinctions are well above the background rate (expected rate of
extinction).

\subsection{Habitat loss: Deforestation}

30\% of forests remaining (1700 baseline)

70\% of forests are within 1km of some human settlement or activity.

\subsection{Why is it important to understand and measure biodiversity?}
\begin{itemize}
	\item Ethical reasons
	\item Functional reasons
	\item Nature's contribution to people (ecosystem services)
\end{itemize}
